% Provenance (never printed): sources/book2/BRIEFS.md, g10-07. Arc and point elasticity, elasticity along a
% line, revenue, income and cross-price elasticities, supply elasticity. Real numbers:
% figdata/grade-10/07-elasticities/ (WDI health spending and GDP per person, 2019); Brenland households
% (invented). Published estimates: ledger g10:gasoline-elasticity, g10:cigarette-elasticity.
\chapter{Elasticities}\label{ch:g10:elasticities}

Line up the countries of the world by income. In 2019 a person in Burundi had an income of about \usd{870} a year, measured
at the prices of a common basket, and the health care spent on him or her came to about \usd{71}. A person in Germany had
\usd{59311} and \usd{6780} of health care. Across the 186 countries of the opening chart, a country with ten times the income
per person spends, on average, about thirteen times as much on health per person. Health care takes a larger share of the income
of rich countries than of poor ones.
% figdata: health (BDI 868, 70.5; DEU 59311, 6779.5); health_fit (n 186, slope 1.12, times 13.2)

``Thirteen times for ten times'' is a ratio of two proportional changes. Economists call such a ratio an elasticity. The grades 6--9
volume measured how strongly the quantity bought responds to the price. This chapter makes that measure precise, shows how it changes
along a demand curve and what it says about the sellers' revenue, and applies the same idea to income, to the prices of other goods
and to supply.

\begin{openingchart}{Health spending per person against GDP per person, 186 countries, 2019, both on times-ten scales, in current
dollars at purchasing-power parity. The line is the straight line, on these scales, that passes closest to the points.}{World Bank,
World Development Indicators: current health expenditure per capita, PPP; GDP per capita, PPP}
\begin{tikzpicture}
\begin{loglogaxis}[omscatter, log basis x=10, log basis y=10, log ticks with fixed point,
  xtickten={3,...,5}, ytickten={1,...,4}, width=0.84\linewidth, height=7.2cm, xmin=600, xmax=150000,
  ymin=10, ymax=20000, xlabel={GDP per person (dollars a year)}, ylabel={health spending per person (dollars a year)}]
\addplot[omDef, only marks, mark=*, mark size=1.2pt, fill opacity=0.7]
  table[x=gdppc, y=health] {figdata/grade-10/07-elasticities/out/health.csv};
\addplot[black, thick, no markers, sharp plot] table[x=gdppc, y=health] {figdata/grade-10/07-elasticities/out/health_line.csv};
\addplot[omThm, only marks, mark=square*, mark size=1.8pt, nodes near coords, point meta=explicit symbolic,
  nodes near coords style={font=\scriptsize, anchor=south east, text=omThm!80!black}]
  table[x=gdppc, y=health, meta=iso3] {figdata/grade-10/07-elasticities/out/health_marked.csv};
\end{loglogaxis}
\end{tikzpicture}
\end{openingchart}

\begin{recall}
From the grades 6--9 volume: the \emph{price elasticity of demand} is the percentage change in the quantity demanded divided by the
percentage change in the price that causes it. Demand is \emph{elastic} when its size is above 1, \emph{inelastic} when below 1, and
\emph{unit-elastic} when it equals 1. A seller's \emph{revenue} is price times quantity. From \cref{ch:g10:demand}: \omterm{def:g10:demand:substitutes}{substitutes},
\omterm{def:g10:demand:substitutes}{complements}, normal and \omterm{def:g10:demand:normal}{inferior goods}. From the grades 6--9 volume
again: a \emph{necessity} takes a smaller share of the budget as income rises, a \emph{luxury} a larger one.
\end{recall}

\begin{notation}
In this chapter $\varepsilon_d$ is a price elasticity of demand, $\varepsilon_s$ a \omterm{def:g10:elasticities:supply}{price elasticity of supply}, $\varepsilon_I$ an
income elasticity and $\varepsilon_{XY}$ the \omterm{def:g10:elasticities:cross-price}{cross-price elasticity} of the demand for $X$ with respect to the price of $Y$. A bar marks a
mean: $\bar Q$ is the mean of two quantities.
\end{notation}

\section{Measuring the response: between two points, and at a point}\label{sec:g10:elasticities:measure}

A bus company raises its fare from 2 to 2.5 dalens, and daily rides fall from \num{1000} to 900, \omterm{def:g10:economic-models:other-things-equal}{other things equal}. In percentages,
the price rose by 25\,\% and rides fell by 10\,\%: the elasticity is $-10/25 = -0.4$. But had the fare fallen from 2.5 to 2, rides would
have risen from 900 to \num{1000}, by 11.1\,\%, while the price fell by 20\,\%: $11.1/(-20) \approx -0.56$. The same two points give two
answers, depending on the starting point. Measuring each change against the mean of the two values removes the problem.

\begin{definition}[Arc elasticity]\label{def:g10:elasticities:arc}
Between two points $(p_1, Q_1)$ and $(p_2, Q_2)$ of a demand curve, the \emph{arc elasticity}\index{arc elasticity} of demand is
\[
  \varepsilon_d = \frac{(Q_2 - Q_1)/\bar Q}{(p_2 - p_1)/\bar p}, \qquad \bar Q = \frac{Q_1 + Q_2}{2}, \quad \bar p = \frac{p_1 + p_2}{2}.
\]
\end{definition}

For the bus: $\bar Q = 950$ and $\bar p = 2.25$, so $\varepsilon_d = (-100/950)/(0.5/2.25) = -0.1053/0.2222 \approx -0.47$, whichever way
the fare moved. Demand for these rides is inelastic.

When the demand curve is known, the response can be measured at a single point, by letting the two points come close together. On a
straight line $Q = a - bp$, the quantity changes by $-b$ for each dalen of price, everywhere.

\begin{definition}[Point elasticity]\label{def:g10:elasticities:point}
At a point $(p, Q)$ of the straight demand line $Q = a - b\,p$, the \emph{point elasticity}\index{point elasticity} of demand is
\[
  \varepsilon_d = -b \, \frac{p}{Q}.
\]
For a curve that is not straight, $-b$ is replaced by the slope $\mathrm{d}Q/\mathrm{d}p$ of the curve at that point.
\end{definition}

\begin{example}[One line, three elasticities]\label{ex:g10:elasticities:line}
For $Q = 120 - 4p$: at $p = 5$, $Q = 100$ and $\varepsilon_d = -4 \times 5/100 = -0.2$; at $p = 15$, $Q = 60$ and $\varepsilon_d = -60/60 =
-1$; at $p = 25$, $Q = 20$ and $\varepsilon_d = -100/20 = -5$. The slope is the same everywhere; the elasticity is not.
\end{example}

Slope and elasticity are different things. The slope says how many units are lost per dalen; it depends on the units in which the good
and the money are measured. The elasticity compares percentages, so it has no unit: it is the same whether bread is counted in loaves or
in tonnes, and prices in dalens or in kests. That is why elasticities can be compared across goods and across countries.

\section{Along a straight demand line}\label{sec:g10:elasticities:line}

\begin{proposition}[Elasticity along a straight line]\label{prop:g10:elasticities:midpoint}
On the demand line $Q = a - b\,p$ (with $a, b > 0$), for $0 < p < a/b$, the size of the \omterm{def:g10:elasticities:point}{point elasticity} is $|\varepsilon_d| = b\,p/(a -
b\,p)$. It rises with the price; it is below 1 when $p < a/(2b)$, equal to 1 at $p = a/(2b)$, the midpoint of the line, and above 1 when $p >
a/(2b)$.
\end{proposition}
\begin{proof}
$|\varepsilon_d| = bp/Q = bp/(a - bp)$. As $p$ rises, the numerator rises and the positive denominator falls, so the ratio rises.
$bp/(a - bp) < 1$ means $bp < a - bp$, that is $p < a/(2b)$; equality holds at $p = a/(2b)$, and the reverse inequality above it. At $p =
a/(2b)$, $Q = a/2$: the point lies halfway between the line's ends, $(0, a/b)$ and $(a, 0)$.
\end{proof}

Near the top of the line the quantity is small, so each lost unit is a large percentage of it, while the price is high, so each dalen is a
small percentage of it: demand is elastic. Near the bottom the reverse holds. Two extreme cases have their own names.

\begin{definition}[Perfectly elastic and perfectly inelastic]\label{def:g10:elasticities:perfectly-elastic}
Demand (or supply) is \emph{perfectly elastic}\index{perfectly elastic} when the curve is horizontal: at one price any quantity is bought,
and at any higher price nothing. It is \emph{perfectly inelastic}\index{perfectly inelastic} when the curve is vertical: the same quantity
is bought whatever the price.
\end{definition}
\phantomsection\label{def:g10:elasticities:perfectly-inelastic}

A single price-taking farmer faces a \omterm{def:g10:elasticities:perfectly-elastic}{perfectly elastic} demand for his wheat: he can sell all he grows at the market price and nothing above
it. A \omterm{def:g10:elasticities:perfectly-elastic}{perfectly inelastic} demand is a limit that real goods only approach: a patient's demand for a life-saving dose, over a short time.

\section{Elasticity and revenue}\label{sec:g10:elasticities:revenue}

When a seller raises the price, each unit sold brings more but fewer units are sold. Which effect wins depends on the elasticity.

\begin{proposition}[Revenue along a straight demand line]\label{prop:g10:elasticities:revenue}
On the demand line $Q = a - b\,p$, revenue $R(p) = p\,(a - b\,p)$ rises with the price where demand is inelastic, falls where it is elastic,
and is largest at the unit-elastic price $p = a/(2b)$, where it equals $a^2/(4b)$.
\end{proposition}
\begin{proof}
$R(p) = ap - bp^2 = -b\,(p - \tfrac{a}{2b})^2 + \tfrac{a^2}{4b}$. This is a parabola open downward with its top at $p = a/(2b)$: it rises for
$p < a/(2b)$ and falls for $p > a/(2b)$. By \cref{prop:g10:elasticities:midpoint}, the first range is where demand is inelastic and the second
where it is elastic. At the top, $R = a^2/(4b)$.
\end{proof}

For $Q = 120 - 4p$, revenue is largest at $p = 15$, with $Q = 60$ and $R = 900$. At $p = 10$ revenue is $10 \times 80 = 800$, and at $p = 20$ it
is $20 \times 40 = 800$ too.

\begin{omfigure}
\begin{tikzpicture}[econ, x=0.075cm, y=0.12cm]
\begin{scope}
\draw[econ axes] (0,33) node[left, font=\small] {$p$} |- (130,0) node[below left=12pt and 0pt, font=\small] {$Q$};
\draw[demand] (0,30) -- (120,0);
\draw[guide] (0,15) -- (60,15) -- (60,0);
\node[eqpoint] at (60,15) {};
\foreach \y in {15,30} \draw (0,\y) -- (-1.5,\y) node[left, font=\footnotesize] {\y};
\foreach \x in {60,120} \draw (\x,0) -- (\x,-0.6) node[below, font=\footnotesize] {\x};
\node[font=\footnotesize, anchor=south west] at (22,25) {elastic, $|\varepsilon_d| > 1$};
\node[font=\footnotesize, anchor=south west] at (62,15.5) {unit-elastic};
\node[font=\footnotesize, anchor=south west] at (88,8) {inelastic, $|\varepsilon_d| < 1$};
\end{scope}
\begin{scope}[yshift=-5.6cm, y=0.0042cm]
\draw[econ axes] (0,1000) node[left, font=\small] {$R$} |- (130,0) node[below left=12pt and 0pt, font=\small] {$Q$};
\draw[curve, domain=0:120, samples=60] plot (\x, {\x*(30 - \x/4)});
\draw[guide] (0,900) -- (60,900) -- (60,0);
\node[eqpoint] at (60,900) {};
\foreach \y in {800,900} \draw (0,\y) -- (-1.5,\y) node[left, font=\footnotesize] {\y};
\draw[guide] (0,800) -- (40,800) -- (40,0);
\draw[guide] (40,800) -- (80,800) -- (80,0);
\foreach \x in {40,60,80,120} \draw (\x,0) -- (\x,-14) node[below, font=\footnotesize] {\x};
\end{scope}
\end{tikzpicture}

{\small Top: the demand line $Q = 120 - 4p$, elastic above its midpoint and inelastic below it. Bottom: revenue $R = pQ$ for each quantity on the
line. Revenue is largest, 900, at the midpoint (60 units at 15 dalens), and equal, 800, at 40 and 80 units.}
\end{omfigure}

The rule gives practical answers. A theatre with empty seats has no reason to cut its prices if demand for its seats is inelastic: it would
fill few more seats and earn less. A government that taxes a good with inelastic demand, such as cigarettes, raises revenue without cutting
consumption much; a tax on a good with elastic demand cuts consumption a lot and raises less. Which goal matters is a separate question; the
elasticity says what each choice does.

\section{Income and other prices}\label{sec:g10:elasticities:income}

The same ratio of percentages measures the response to any cause.

\begin{definition}[Income elasticity of demand]\label{def:g10:elasticities:income}
The \emph{income elasticity of demand}\index{income elasticity of demand} for a good is the percentage change in the quantity demanded (or in
spending on it, prices held fixed) divided by the percentage change in income that causes it, prices held fixed:
$\varepsilon_I = (\Delta Q/\bar Q)/(\Delta I/\bar I)$.
\end{definition}

A \omterm{def:g10:demand:normal}{normal good} has $\varepsilon_I > 0$ and an \omterm{def:g10:demand:normal}{inferior good} $\varepsilon_I < 0$. In these terms, a necessity, defined in the grades 6--9 volume,
has $0 < \varepsilon_I < 1$: its share of the budget falls as income rises; a luxury has $\varepsilon_I > 1$: its share rises.

\begin{omchart}{Brenland: food spending against disposable income, \num{2000} households, one year, in dalens. The share of food falls from
the poorest fifth of households to the richest.}{Brenland household survey (invented data)}
\begin{tikzpicture}
\begin{axis}[omscatter, width=0.8\linewidth, height=6cm, xmin=0, xmax=90000, ymin=0, ymax=10000,
  xlabel={disposable income (dalens a year)}, ylabel={food spending (dalens a year)},
  scaled ticks=false, xticklabel style={/pgf/number format/1000 sep={\,}}, yticklabel style={/pgf/number format/1000 sep={\,}}]
\addplot[omDef, only marks, mark=*, mark size=0.6pt, fill opacity=0.4, draw opacity=0.4]
  table[x=disposable, y=food] {figdata/grade-10/07-elasticities/out/food.csv};
\end{axis}
\end{tikzpicture}
\end{omchart}

In the Brenland survey, the poorest fifth of households had a mean disposable income of \dal{13870} and spent \dal{3940} on food, 28.4\,\% of their
income; the richest fifth had \dal{41790} and spent \dal{5340}, 12.8\,\%. Between the two, the arc income elasticity of food spending is
\[
  \varepsilon_I = \frac{1400/4640}{27920/27830} \approx \frac{0.302}{1.003} \approx 0.30 .
\]
A household three times richer spends about a third more on food. Food is a necessity, and Brenland behaves like the real household surveys
described in the grades 6--9 volume.
% figdata: food_fit (inc_low 13870, inc_high 41790, food_low 3940, food_high 5340, shares 28.4 and 12.8, arc 0.30)

The opening chart reads the same way across countries: a slope of 1.12 on the two times-ten scales says that 1\,\% more income per person goes
with 1.12\,\% more health spending per person, so health care behaves there like a luxury. That reading needs care. The chart compares
different countries, not one country whose income grew: richer countries also have older populations, dearer doctors and different
systems of paying for care. It is a pattern across countries, not a demand curve.
% figdata: health_fit (slope 1.12)

\begin{definition}[Cross-price elasticity]\label{def:g10:elasticities:cross-price}
The \emph{cross-price elasticity}\index{cross-price elasticity} of the demand for good $X$ with respect to the price of good $Y$ is the percentage
change in the quantity of $X$ demanded divided by the percentage change in the price of $Y$: $\varepsilon_{XY} = (\Delta Q_X/\bar Q_X)/(\Delta
p_Y/\bar p_Y)$.
\end{definition}

It is positive for \omterm{def:g10:demand:substitutes}{substitutes} (a dearer bus sends people to the tram) and negative for \omterm{def:g10:demand:substitutes}{complements} (dearer petrol means fewer car journeys).
Close to zero, the two goods hardly affect each other.

\section{Supply, time and what makes demand elastic}\label{sec:g10:elasticities:supply}

\begin{definition}[Price elasticity of supply]\label{def:g10:elasticities:supply}
The \emph{price elasticity of supply}\index{price elasticity of supply} is the percentage change in the quantity supplied divided by the percentage
change in the price that causes it: $\varepsilon_s = (\Delta Q^s/\bar Q^s)/(\Delta p/\bar p)$, or, on a supply line $Q = c + d\,p$, $\varepsilon_s =
d\,p/Q$ at a point.
\end{definition}

It is positive for an upward-sloping supply curve. \Cref{ch:g10:supply} showed why it grows with time: the long-run supply of crude oil or of soya
responds far more than the short-run one. The same holds for demand. Three things make demand for a good more elastic.

\begin{itemize}
\item \textbf{Close \omterm{def:g10:demand:substitutes}{substitutes}.} The demand for one brand of petrol is very elastic, because the station across the road sells nearly the same
  thing; the demand for petrol as a whole is not.
\item \textbf{A large share of the budget.} A 10\,\% rise in the price of salt is hardly noticed; a 10\,\% rise in rents changes where people live.
\item \textbf{Time.} In a week, drivers can only drive a little less. In ten years they can buy smaller cars, move closer to work, or take a job
  nearer home.
\end{itemize}

\begin{measure}[How an elasticity is estimated: comparing like with like]
Joining two observed points of price and quantity gives an elasticity only if nothing else moved demand between them. \Cref{ch:g10:demand} showed
that the yearly points of the oil market mix moves along the demand curve with shifts of it. Careful studies look for price changes that came from
the supply side and say nothing about buyers: a tax change, a frost, a strike at a refinery. They compare the quantities before and after with the
quantities of similar buyers whom the change did not touch. A survey of many studies of petrol, of every design, gives a size near 0.09 within a year and
near 0.31 in the long run, once it corrects for the habit of publishing large estimates; for cigarettes in the United States, 0.3 to 0.6 for adults and 0.6 to 0.7 for young people. The estimates differ from
study to study, and reporting them as a range is honest.
\end{measure}
% ledger: g10:gasoline-elasticity (-0.09 short run, -0.31 long run); g10:cigarette-elasticity (-0.3 to -0.6 general, -0.6 to -0.7 youth)

\begin{omfigure}
\small
\begin{tabular}{l c c}
\toprule
good & price elasticity of demand & time span \\
\midrule
car fuel (many countries) & $-0.09$ & short run \\
car fuel (many countries) & $-0.31$ & long run \\
cigarettes (United States, all) & $-0.3$ to $-0.6$ & --- \\
cigarettes (United States, young people) & $-0.6$ to $-0.7$ & --- \\
\bottomrule
\end{tabular}

{\small Published estimates of price elasticities of demand. Sources: a survey of studies of car-fuel demand (Havranek, Irsova and Janda 2012);
a textbook summary of cigarette studies (Perloff 2008).}
\end{omfigure}
% ledger: g10:gasoline-elasticity; g10:cigarette-elasticity

\section{Exercises}

\begin{exercise}[$\star$]\label{exo:g10:elasticities:1}
When the price of a ticket rises from 4 to 5 dalens, the number sold falls from 600 to 500. Compute the \omterm{def:g10:elasticities:arc}{arc elasticity} of demand. Is demand elastic
or inelastic in this range?
\end{exercise}

\begin{exercise}[$\star$]\label{exo:g10:elasticities:2}
When the price of a sandwich falls from 10 to 8 dalens, sales rise from 200 to 260 a day. Compute the \omterm{def:g10:elasticities:arc}{arc elasticity}. Does the shop's revenue rise or
fall?
\end{exercise}

\begin{exercise}[$\star$]\label{exo:g10:elasticities:3}
Read the labelled points of the opening chart: the United States (GDP per person \usd{65228}, health spending \usd{10550}) and India (\usd{6945} and
\usd{211}). How many times higher are US income and US health spending per person? Is the United States above or below the line?
\end{exercise}

\begin{exercise}[$\star$]\label{exo:g10:elasticities:4}
Give the sign of the \omterm{def:g10:elasticities:cross-price}{cross-price elasticity} of the demand for the first good with respect to the price of the second: (a) tea and coffee; (b) cars and
petrol; (c) bread and butter; (d) bus tickets and train tickets between the same two towns.
\end{exercise}

\begin{exercise}[$\star$]\label{exo:g10:elasticities:5}
Classify a demand with each price elasticity: $-0.2$; $-1$; $-3$; $0$; a horizontal demand curve.
\end{exercise}

\begin{exercise}[$\star\star$]\label{exo:g10:elasticities:6}
For the demand line $Q = 200 - 5p$, compute the \omterm{def:g10:elasticities:point}{point elasticity} at $p = 10$, $p = 20$ and $p = 30$. At which price is demand unit-elastic?
\end{exercise}

\begin{exercise}[$\star\star$]\label{exo:g10:elasticities:7}
For the supply line $Q = -20 + 4p$, compute the \omterm{def:g10:elasticities:point}{point elasticity} of supply at $p = 10$ and at $p = 20$.
\end{exercise}

\begin{exercise}[$\star\star$]\label{exo:g10:elasticities:8}
For the demand line $Q = 90 - 3p$, find the price that gives the largest revenue, the quantity sold and the revenue. Check that demand is
unit-elastic there.
\end{exercise}

\begin{exercise}[$\star\star$]\label{exo:g10:elasticities:9}
Using the fifth means of the Brenland survey given in the text, compute the food share of the poorest and of the richest fifth, and the arc income
elasticity of food spending between them. What would the elasticity have been if the richest fifth had spent the same share of its income on food as
the poorest?
\end{exercise}

\begin{exercise}[$\star\star$]\label{exo:g10:elasticities:10}
A museum sells \num{5000} tickets a month at 10 dalens. Its demand has a price elasticity of $-0.4$. It raises the price by 10\,\%. Treating the
elasticity as the ratio of the two percentage changes, estimate the new number of tickets and the change in revenue.
\end{exercise}

\begin{exercise}[$\star\star$]\label{exo:g10:elasticities:11}
Use the table of published estimates. (a) A lasting 10\,\% rise in the price of car fuel: by about how much does fuel use fall within a year, and in the
long run? (b) A tax raises the price of cigarettes by 20\,\%: by how much does adult smoking fall, at the two ends of the range, and young people's?
\end{exercise}

\begin{exercise}[$\star\star\star$]\label{exo:g10:elasticities:12}
A study claims that the demand for a good is inelastic. Describe an observation that would contradict this claim, and say what you would check about how
the observation was made.
\end{exercise}

\begin{exercise}[$\star\star\star$]\label{exo:g10:elasticities:13}
On the line $Q = a - b\,p$, let $P = a/b$ be the price at which nobody buys. Prove that the \omterm{def:g10:elasticities:point}{point elasticity} at a price $p$ is $\varepsilon_d = -p/(P - p)$.
Deduce, without computing $Q$, the price at which demand is unit-elastic and the elasticity at a price of $3P/4$.
\end{exercise}

\begin{exercise}[$\star\star\star$]\label{exo:g10:elasticities:14}
The demand for a good is $Q = 1200/p$. Show that buyers spend the same amount whatever the price, and that the \omterm{def:g10:elasticities:arc}{arc elasticity} between $p = 2$ and $p = 3$ is
exactly $-1$. Describe in words a demand curve whose elasticity is the same at every point, and say why it is not a straight line.
\end{exercise}

\begin{exercise}[$\star\star\star$]\label{exo:g10:elasticities:15}
Explain, with car fuel as the example, why long-run price elasticities of demand are larger than short-run ones. Then explain, with
\cref{ex:g10:supply:two-curves}, why the same holds for supply.
\end{exercise}

\section{Problem: The Toll Bridge}

\begin{problem}[{Weekend problem --- a toll bridge, a revenue-maximising toll and a new ferry}]
\label{pb:g10:elasticities:1}
A river town in Brenland owns a toll bridge. In one year the toll was 4 dalens and there were \num{9000} crossings a day; the next year the toll was 5
dalens and there were \num{8000}. Nothing else that matters to travellers changed. The town treats demand as a straight line through these two points.

\medskip
\noindent\textbf{Part I --- The demand for crossings.}
\begin{enumerate}
\item Find the demand line $Q = a - b\,p$ through the two points.
\item At what toll would nobody cross? How many crossings would there be with no toll?
\item Compute the \omterm{def:g10:elasticities:point}{point elasticity} of demand at a toll of 5 dalens.
\item Compute the \omterm{def:g10:elasticities:arc}{arc elasticity} between the two years.
\item Compute the town's revenue in each year. Is this consistent with the elasticity?
\end{enumerate}

\medskip
\noindent\textbf{Part II --- The revenue-maximising toll.}
\begin{enumerate}[resume]
\item Write the revenue $R(p)$ as a function of the toll.
\item Find the toll that brings the largest revenue, and the number of crossings at that toll.
\item Compute that revenue and the elasticity at that toll.
\item Check: compute the revenue at 6 and at 7 dalens.
\item Compute the revenue at 8 dalens. Why does a higher toll eventually bring less?
\end{enumerate}

\medskip
\noindent\textbf{Part III --- The ferry.}
\begin{enumerate}[resume]
\item A ferry starts to cross the river, with a fare of 3 dalens, and the demand for the bridge becomes $Q = 12000 - 1200\,p$. How many crossings does the
  bridge lose at a toll of 5 dalens?
\item Compute the \omterm{def:g10:elasticities:point}{point elasticity} at 5 dalens now.
\item Find the revenue-maximising toll and the revenue it brings.
\item If the town had already moved the toll to the value found in question 7, what would its revenue be with the ferry? Compare with question 13.
\item Explain why the ferry made the demand for the bridge more elastic.
\end{enumerate}

\medskip
\noindent\textbf{Part IV --- Other prices and incomes.}
\begin{enumerate}[resume]
\item With a ferry fare of $f$ dalens, the demand for the bridge is $Q = 10800 - 1200\,p + 400\,f$. Check that this gives the demand of question 11 when $f
  = 3$.
\item Compute the \omterm{def:g10:elasticities:cross-price}{cross-price elasticity} of bridge crossings with respect to the ferry fare, at $p = 5$ and $f = 3$. What does its sign say?
\item The ferry raises its fare to 4 dalens. How many crossings, and what revenue, at a toll of 5?
\item The incomes of the town's travellers rise by 10\,\%, and at every toll and fare the number of crossings rises by 3\,\%. Estimate the income elasticity
  of crossings. Is a bridge crossing a \omterm{def:g10:demand:normal}{normal good} here, and a necessity or a luxury?
\item State the result of the problem: the revenue-maximising toll and the revenue it brings, before and after the ferry.
\end{enumerate}
\end{problem}
